\section{Discussions}
Despite the wide usage in DG benchmarks, Adam is hardly the optimal choice of default optimizer. We show that the best optimizer for different DG methods varies across different datasets. Choosing the right optimizer can help exceed the current benchmarks' limitations. We propose a unified optimization method, flatness-Aware minimization for Domain generalization (FAD), to optimize zeroth-order and first-order flatness efficiently to address the problem. We further show that FAD controls the OOD generalization error and convergences fast. Compared with current optimizers, flatness-aware optimization approaches, including SAM, GAM, and FAD, leads to flatter minima and better generalization. FAD outperforms its counterparts significantly on most DG datasets. 

FAD still has the following limitations which could lead to potential future work. First, the term $\sup_{\bftheta_1, \bftheta_2 \in \caltheta}\left|\calL_{\train}(\bftheta_1, \bftheta_2) - \calL_{\test}(\bftheta_1, \bftheta_2)\right|$ in Equation \eqref{eq:main-bound} indicates the discrepancy between training and test data. Without knowledge of test distribution and extra assumptions, the term cannot be optimized and leads to the optimization of the flatness term. Proper assumptions on the connection between training and test data or the geometry of data may help to learn the connection between flatness and the discrepancy, and thus control the discrepancy for better generalization. Second, We show that Adam is hardly an optimal choice of default optimizer in DG. However, selecting a proper optimizer for each DG method leads to massive overhead. Designing a better DG protocol or learning an indicator for the fitness of optimizers and methods can significantly contribute.  